\begin{proposition*}[Relations on one vector; Midterm 2 Review, Problem 6]
For $v\in V$ there is a unique monic polynomial $p_v$ of least degree with $p_v(T)v=0$. For every polynomial $q$,
\[q(T)v=0\iff p_v\mid q;\qquad p_v\mid m_T.\]
For $v=0$, $p_v=1$.
\end{proposition*}
\begin{prob}\begin{proof}\begin{steps}
\ii For $v\ne0$, the first dependence among $v,Tv,\ldots$ gives a relation; divide by its last nonzero coefficient to make it monic. Two least-degree monic choices have a smaller-degree difference with zero value on $v$, so they agree.
\ii Divide $q=p_vh+r$, with $r=0$ or $\deg r<\deg p_v$. If $q(T)v=0$, commutation gives $r(T)v=q(T)v-h(T)p_v(T)v=0$. Minimality forces $r=0$. Conversely, a multiple of $p_v$ has zero value on $v$ by the same product equation.
\ii Apply this to $q=m_T$. For $v=0$, every polynomial gives zero, and the unique monic degree-zero choice is $1$.
\end{steps}\end{proof}\end{prob}
\begin{prob}\begin{definition*}[Least common multiple]
The monic least common multiple $\operatorname{lcm}(p_1,\ldots,p_r)$ of nonzero polynomials is a monic common multiple that divides every common multiple.
\end{definition*}
\begin{steps}\ii The product is a common multiple, so choose a monic common multiple $L$ of least degree.
\ii Divide any common multiple $q=Lh+r$. Every $p_j$ divides $q$ and $Lh$, hence divides $r$. A nonzero $r$ would be a common multiple of smaller degree; thus $r=0$ and $L\mid q$.
\ii Two such monic choices divide each other and therefore agree. This proves existence and uniqueness.
\end{steps}\end{prob}
\begin{proposition*}[Combine relations on a basis]
For any basis $v_1,\ldots,v_n$,
\[m_T=\operatorname{lcm}(p_{v_1},\ldots,p_{v_n}).\]
The empty least common multiple is $1$ on $V=\{0\}$.
\end{proposition*}
\begin{prob}\begin{proof}\begin{steps}
\ii A linear map $q(T)$ is zero iff it is zero on every basis vector: expand $q(T)(\sum_j a_jv_j)=\sum_j a_jq(T)v_j$.
\ii By the vector relation criterion, this holds iff every $p_{v_j}$ divides $q$. Their least common multiple and $m_T$ therefore have exactly the same polynomial multiples; each divides the other and both are monic, so they agree.
\end{steps}\end{proof}\end{prob}
